% ============================================================
% M3 S2 导学件·波1 片件 —— 课时05 两点式与一般式（2.2.2下）
% 生成：M3 成卷轮 S2 波1臂2｜2026-09-14｜编译 cwd＝本目录（中文路径禁 \input 绝对路径）
%   xelatex main-true.tex ×2 → 含详解印本；xelatex main-false.tex ×2 → 纯题版（单源双本）
% 输入链（全只读）：题面库 课时05-两点式与一般式.md（题面逐字装配；【注】行＝装配层注记不印，
%   题首「（提示：…）」行＝印面照录）＋ -答案侧.md（值逐字回嵌）
%   ＋ 定稿/批A-课时05-两点式与一般式.md（详解回嵌源）；
%   toolchain＝qp-m3.sty（钉后版 md5 7c3930362be8a0a2bdf21bbf8ac16573，本地挂载副本，破损即停）。
% 母版体例：承 成卷/导学件/课时01-坐标法/母版体例说明.md（骨架/宏用法/门谱跑法原样照抄）。
% 键账：件内 ansblock 键集 ≡ 题面库片键集 ≡ manifest 键序（24 键，零缺漏零多余）；
%   预习位零键零答案（口令④前口径）；印面号 1—24＝探究 1—19＋课堂评价 20—24，
%   键↔号映射见 值台账-课时05.json。
% 片特记（跳号对号档，题面侧§片特记在案）：原练习 E1/E4/E8/E12/E13/E14 六席已前移课时04
%   （＝04-E11~E16，本片留空号非键）；原拓展 T2/T3/T4/T5/T9/T10 已回编练习 E17~E22；
%   05-T1 为占位行不计键不计题；现拓展实题＝T6/T7/T8。装配层不补号、不重编号。
% 括线判据（承重墙禁放宽，toolchain钉档§三.1）：估高＞8 行（chars/23 栏宽口径）→ 括线模，
%   逐键读数入值台账；其余灰底。本片括线键（12）：E3、E15、E16、E17、E18、E19、E20、E21、E22、
%   T6、T7、T8（估高9/11/16/10/11/16/26/11/12/13/14/21 行）；灰底 12 键（E2/E5/E6/E7/E9/E10＋G1~G5）。
% 自撰文注：素养小结/目标/预习条目为本片自撰（非源件逐字），其中探究点三③收束句避用「答案」一词
%   （false 档泄答词门禁该字出现在印面），已改「源卷结果」。
% 门谱读数：见 过程件 工作区/_tmpM3S2波1臂2_0914/（编译 log、门谱输出、目验 PNG 双档）。
% ============================================================
\documentclass[fontset=none]{ctexart}
\usepackage{qp-m3}
% —— 印面逐字符号路由补钉（探针实证 0914：书宋缺 U+2212，▱ 诸字库皆缺→NSC 子块；
%     禁改 sty 本体保 md5；无 ▱/− 用件的片可整块照抄，零副作用） ——
\xeCJKDeclareCharClass{Default}{"2212}%
\setCJKfamilyfont{nscblk}[Path=C:/提示词/工作区/字替对照-0909/variantF/fonts/,BoldFont=NSC-w600.ttf]{NSC-w500.ttf}%
\xeCJKDeclareSubCJKBlock{nscblk}{"25B1}%
\newcommand{\pxparallelogram}{{\CJKfamily{nscblk}▱}}%
% —— 断行微调（纯题档/详解档三零门：CJK 胶收缩 0.01→0.05em；应急伸展曾试 12em/2em，
%     2026-09-14 实测撤行回落 sty 默认 1em 三零仍守，故不设该行——补丁最少化；
%     只动伸缩分量不动自然宽度，版面纹理零漂；Underfull 治理读数见过程件） ——
\xeCJKsetup{CJKglue={\hskip 0.10em plus 0.08em minus 0.05em}}%
% —— 页脚件名（qp-headfoot 复用入口） ——
\renewcommand{\qpzhangming}{第二章\quad 平面解析几何}%
\renewcommand{\qpjianming}{导学件}%

\begin{document}

\zhangtitle{第二章\quad 平面解析几何}
\jietitle{2.2.2 直线的两点式与一般式}
\mubiaoline
\mubiaomu{1}{掌握直线方程的两点式、截距式与一般式，理解各种形式的适用条件与局限，能熟练完成五种形式间的互化；}
\mubiaomu{2}{能由二元一次方程的系数判定直线的位置特征（过原点、与坐标轴平行或垂直、截距情形），并会求直线与两坐标轴围成图形的面积；}
\mubiaomu{3}{会用含参直线系整理定点的方法处理恒过定点问题，体会“形式选择服务于条件”与“数形结合、以算代证”的思想.}

\vspace{1.7mm}
\begin{multicols}{2}
\emergencystretch=1em

% ============================================================
% 课前预习（知识点导学填空；零键零答案——预习键位按检查口令④裁定前口径，
% \kongbai 空线＋\zhentib 题面形，两档等形不泄答）
% ============================================================
\huaxing{课}{前}{预}{习}{知识导学\quad 素养初识}

\zsd{一}{直线的两点式}

\tiaomu{1}{经过两点\(P_1(x_1,y_1)\)、\(P_2(x_2,y_2)\)（\(x_1\neq x_2\)且\(y_1\neq y_2\)）的直线方程为\((y-y_1)/(y_2-y_1)=\)\kongbai{}，此方程叫做直线方程的两点式．\zhuzhu{两点式表示不了与坐标轴平行（或重合）的直线：\(x_1=x_2\) 时直线为\(x=x_1\)，\(y_1=y_2\) 时直线为\(y=y_1\)，须单列；写成\((y_2-y_1)(x-x_1)=(x_2-x_1)(y-y_1)\)的整式形则对任意两点都适用.}}

\zsd{二}{直线的截距式}

\tiaomu{1}{若直线\(l\)与\(x\)轴、\(y\)轴的交点分别为\((a,0)\)、\((0,b)\)（\(ab\neq0\)），则\(l\)的方程为\(x/a+y/b=\)\kongbai{}，此方程叫做直线方程的截距式，其中\(a\)、\(b\)分别是直线在\(x\)轴、\(y\)轴上的截距．\zhuzhu{截距式要求两截距都存在且不为零：过原点的直线、与坐标轴平行的直线都没有截距式；求与两轴围成面积的题常先设截距式.}}

\zsd{三}{直线的一般式}

\tiaomu{1}{关于\(x\)、\(y\)的二元一次方程\(Ax+By+C=\)\kongbai{}（其中\(A\)、\(B\)不同时为0）叫做直线的一般式方程．\zhuzhu{五种形式中只有一般式能表示所有直线；点斜式、斜截式不含竖直线，两点式、截距式另有局限，故“设方程”时形式选择要按条件定，求出后一般化简为一般式.}}

\tiaomu{2}{直线\(Ax+By+C=0\)（\(B\neq0\)）的斜率为\kongbai{}，在\(y\)轴上的截距为\kongbai{}．\zhuzhu{化斜截式\(y=-(A/B)x-(C/B)\)即得；\(B=0\)时直线与\(x\)轴垂直，无斜率.}}

\zhenhead{判断正误(正确的打\gou{},错误的打\cha{})}

\zhentib{(1)}{直线\(x/a+y/b=1\)（\(ab\neq0\)）在\(x\)轴、\(y\)轴上的截距分别为\(a\)、\(b\)．}

\zhentib{(2)}{方程\((y_2-y_1)(x-x_1)=(x_2-x_1)(y-y_1)\)只能表示与两坐标轴都不平行的直线．}

\liubai[9mm]{此处书写}

% ============================================================
% 课中探究（考点探究〔例+变式〕＝练/拓槽 21 键；题后紧跟答案＝ansblock 内嵌，
% 值照答案侧逐字，详解承批A 回嵌＋印面清洗注记见值台账）
% ============================================================
\huaxing{课}{中}{探}{究}{考点探究\quad 素养小结}

% ---------- 探究点一 两点式及其适用条件 ----------
\tjdnr{一}{两点式及其适用条件}{例\textbf{1}}{中档(知识点一)}{若直线\(l\)过点\((-1,-1)\)和\((2,5)\)，且点\((1009,b)\)在直线\(l\)上，则\(b\)的值为\nobreak(\kongwei)}

\bindopt {\fontsize{10.5pt}{\qpTJgLeadLiti}\selectfont \makebox[\dimexpr0.25\linewidth-0.5em\relax][l]{A．2019}\makebox[\dimexpr0.25\linewidth-0.5em\relax][l]{B．2018}\makebox[\dimexpr0.25\linewidth-0.5em\relax][l]{C．2017}\makebox[\dimexpr0.25\linewidth-0.5em\relax][l]{D．2016}\par}

\begin{ansblock}[2章-练-课时05-E2]
% ans:2章-练-课时05-E2
\ansitem{1}{A}
\ansnote{详解}{两点式：\((y+1)/(5+1)=(x+1)/(2+1)\)，即\((y+1)/6=(x+1)/3\)，化简得\(y=2x+1\)．\par\noindent 代入\((1009,b)\)：\(b=2\times 1009+1=2019\)．选：A．}
\end{ansblock}

\liB{变式\textbf{2}}{简单(知识点一)}{}{过\(A(x_1,y_1)\)和\(B(x_2,y_2)\)两点的直线方程是\nobreak(\kongwei)}

\lxopt{A．\((y-y_1)/(y_2-y_1)=(x-x_1)/(x_2-x_1)\)}

\lxopt{B．\((y-y_1)/(x-x_1)=(y_2-y_1)/(x_2-x_1)\)}

\lxopt{C．\((y_2-y_1)(x-x_1)-(x_2-x_1)(y-y_1)=0\)}

\lxopt{D．\((x_2-x_1)(x-x_1)-(y_2-y_1)(y-y_1)=0\)}

{\ansblockgrayfalse
\begin{ansblock}[2章-练-课时05-E3]
% ans:2章-练-课时05-E3
\ansitem{2}{C}
\ansnote{详解}{A、B要求\(x_1\neq x_2\)且\(y_1\neq y_2\)（或至少分母不为0），不能表示与坐标轴平行（或重合）的直线；\par\noindent C式即\((y_2-y_1)(x-x_1)=(x_2-x_1)(y-y_1)\)：当\(x_1\neq x_2\)时化为点斜式，当\(x_1=x_2\)时得\(x=x_1\)（竖直线）也成立，\(y_1=y_2\)时得\(y=y_1\)同样成立，对任意两点均恰当表示直线；D是把两坐标“同向”组合，一般不表示过A、B的直线（它表示到A、B横、纵坐标差相等点的轨迹型式子）．选：C．}
\end{ansblock}}

\liB{变式\textbf{3}}{简单(知识点一)}{}{已知点\(A(3,2)\)，\(B(-1,4)\)，则过点\(C(2,5)\)且过线段AB的中点的直线方程为\kongbai{}．}
\xiexwei{8mm}

\begin{ansblock}[2章-练-课时05-E11]
% ans:2章-练-课时05-E11
\ansitem{3}{2x−y＋1＝0}
\ansnote{详解}{AB中点\(D((3-1)/2,(2+4)/2)=(1,3)\)．\(C\)、\(D\)横坐标不等（\(2\neq1\)）、纵坐标不等（\(5\neq3\)），用两点式：\((y-5)/(3-5)=(x-2)/(1-2)\)，即\((y-5)/(-2)=(x-2)/(-1)\)，\par\noindent 去分母：\(y-5=2(x-2)\)，化简得\(2x-y+1=0\)．}
\end{ansblock}

\liB{变式\textbf{4}}{中档(知识点三)}{}{判断\((y-2)/(x-3)=1\)是否是某条直线的方程．}
\xiexwei{10mm}

{\ansblockgrayfalse
\begin{ansblock}[2章-练-课时05-E15]
% ans:2章-练-课时05-E15
\ansitem{4}{不是（见详解）}
\ansnote{详解}{方程\((y-2)/(x-3)=1\)要求\(x\neq3\)，其解集对应的点集是直线\(y-2=x-3\)（即\(y=x-1\)）上去掉点\((3,2)\)后的全体；\par\noindent 而“直线\(y=x-1\)的方程”应使以该直线上的点的坐标为解、且以方程的解为坐标的点都在这条直线上——点\((3,2)\)在直线\(y=x-1\)上却不是该方程的解，两集合不同．\par\noindent 所以\((y-2)/(x-3)=1\)不是（这条）直线的方程，它表示去掉一个点的“直线型”曲线；应改写成\((y-2)=(x-3)\)即\(x-y-1=0\)．}
\end{ansblock}}

\xiaojie{①识别：以由两点求方程、辨析各式适用条件、判定方程与直线是否同解为目标的题，判归本题型；先查两点横、纵坐标是否相等再选形式.}
\bindp {\kaishu ②操作：两点式先验\(x_1\neq x_2\)且\(y_1\neq y_2\)，不满足者直接写\(x=x_1\)或\(y=y_1\)；求中点、分点先用中点公式坐标代入.}
\bindp {\kaishu ③收束：方程与直线一一对应要求解集与点集完全相同，去分母时若丢点须回代检验（分母含未知数的变形最易失解）.}


% ---------- 探究点二 直线的一般式与系数关系 ----------
\tjdnr{二}{直线的一般式与系数关系}{例\textbf{5}}{中档(知识点三)}{直线的方程\(Ax+By+C=0\)（\(A^{2}+B^{2}\neq0\)）中的\(A\)，\(B\)，\(C\)满足什么条件时，直线分别具有如下性质？}

\bindp (1) 过坐标原点； (2) 与两条坐标轴都相交； (3) 与\(x\)轴无交点；

\bindp (4) 与\(y\)轴无交点； (5) 与\(x\)轴垂直； (6) 与\(y\)轴垂直．
\xiexwei{12mm}

{\ansblockgrayfalse
\begin{ansblock}[2章-练-课时05-E16]
% ans:2章-练-课时05-E16
\ansitem{5}{(1) C＝0；(2) A≠0且B≠0；(3) A＝0且C≠0；(4) B＝0且C≠0；(5) B＝0（且A≠0）；(6) A＝0（且B≠0）}
\ansnote{详解}{(1) 原点坐标满足方程：\(C=0\)（\(A\)、\(B\)不同时为0由前提保证）；\par\noindent (2) 与两坐标轴都相交\(\iff\)既不与\(x\)轴平行也不与\(y\)轴平行\(\iff A\neq0\)且\(B\neq0\)；\par\noindent (3) 与\(x\)轴无交点\(\iff\)\(l\)为水平线\(y=m\)（\(A=0\)，\(B\neq0\)）且\(m=-C/B\neq0\)，即\(A=0\)且\(C\neq0\)；\par\noindent (4) 与\(y\)轴无交点\(\iff\)\(l\)为竖直线\(x=n\)（\(B=0\)，\(A\neq0\)）且\(n=-C/A\neq0\)，即\(B=0\)且\(C\neq0\)；\par\noindent (5) 与\(x\)轴垂直\(\iff\)\(l\)为竖直线\(x=n\)（\(n\)可为0）：\(B=0\)且\(A\neq0\)；\par\noindent (6) 与\(y\)轴垂直\(\iff\)\(l\)为水平线\(y=m\)（\(m\)可为0）：\(A=0\)且\(B\neq0\)．\par\noindent （说明：(3)(4)是(5)(6)中去掉与坐标轴重合情形的细分——(3)即(6)中\(C\neq0\)（排除\(y=0\)），(4)即(5)中\(C\neq0\)（排除\(x=0\)）．）}
\end{ansblock}}

\liB{变式\textbf{6}}{中档(知识点三)}{}{若方程\((2m^{2}+m-3)x+(m^{2}-m)y-4m+1=0\)表示一条直线，则实数\(m\)满足\nobreak(\kongwei)}

\lxopt{A．\(m\neq0\)}

\lxopt{B．\(m\neq-3/2\)}

\lxopt{C．\(m\neq1\)}

\lxopt{D．\(m\neq1\)且\(m\neq-3/2\)且\(m\neq0\)}

\begin{ansblock}[2章-练-课时05-E5]
% ans:2章-练-课时05-E5
\ansitem{6}{C}
\ansnote{详解}{二元一次方程表示直线\(\iff\)x、y系数不全为0．令\(2m^{2}+m-3=(2m+3)(m-1)=0\)且\(m^{2}-m=m(m-1)=0\)同时成立，公共解只有\(m=1\)；\par\noindent 故方程表示直线\(\iff m\neq1\)（\(m=-3/2\)时仅\(x\)系数为0、\(y\)系数\(15/4\neq0\)，仍表示直线；\(m=0\)时\(x\)系数\(-3\neq0\)）．选：C．}
\end{ansblock}

\liB{变式\textbf{7}}{简单(知识点三)}{}{若点\(P(x_0,y_0)\)是直线\(l\)：\(Ax+By+C=0\)外一点，则方程\(Ax+By+C+(Ax_0+By_0+C)=0\)表示\nobreak(\kongwei)}

\lxopt{A．过点\(P\)且与\(l\)平行的直线}

\lxopt{B．过点\(P\)且与\(l\)垂直的直线}

\lxopt{C．不过点\(P\)且与\(l\)平行的直线}

\lxopt{D．不过点\(P\)且与\(l\)垂直的直线}

\begin{ansblock}[2章-练-课时05-E6]
% ans:2章-练-课时05-E6
\ansitem{7}{C}
\ansnote{详解}{\(P\)在\(l\)外\(\Rightarrow Ax_0+By_0+C\neq0\)．新方程即\(Ax+By+2C+Ax_0+By_0=0\)，\(x\)、\(y\)系数与\(l\)完全相同而常数项不同，故新直线与\(l\)平行（不重合）；\par\noindent 将\(P\)代入新方程左边得\((Ax_0+By_0+C)+(Ax_0+By_0+C)=2(Ax_0+By_0+C)\neq0\)，故新直线不过\(P\)．选：C．}
\end{ansblock}

\liB{变式\textbf{8}}{简单(知识点一)}{}{已知直线\(a_1x+b_1y+1=0\)和直线\(a_2x+b_2y+1=0\)都过点\(A(2,1)\)，则过点\(P_1(a_1,b_1)\)和点\(P_2(a_2,b_2)\)的直线方程是\nobreak(\kongwei)}

\bindopt {\fontsize{10.5pt}{\qpTJgLeadLiti}\selectfont \makebox[\dimexpr0.5\linewidth-1em\relax][l]{A．\(2x+y+1=0\)}\makebox[\dimexpr0.5\linewidth-1em\relax][l]{B．\(2x-y+1=0\)}\par}

\bindopt {\fontsize{10.5pt}{\qpTJgLeadLiti}\selectfont \makebox[\dimexpr0.5\linewidth-1em\relax][l]{C．\(2x+y-1=0\)}\makebox[\dimexpr0.5\linewidth-1em\relax][l]{D．\(x+2y+1=0\)}\par}

\begin{ansblock}[2章-练-课时05-E7]
% ans:2章-练-课时05-E7
\ansitem{8}{A}
\ansnote{详解}{把\(A(2,1)\)代入两线方程：\(2a_1+b_1+1=0\)且\(2a_2+b_2+1=0\)，\par\noindent 即\(P_1\)、\(P_2\)两点都满足方程\(2x+y+1=0\)．又两已知直线不同（否则\(a_1:a_2=b_1:b_2=1:1\)，题中\(P_1\)、\(P_2\)为两点确定直线），所以\(2x+y+1=0\)即所求直线．选：A．}
\end{ansblock}

\liB{变式\textbf{9}}{简单(知识点三)}{}{点\(M(x_0,y_0)\)是直线\(Ax+By+C=0\)上的点，则直线方程可表示为\nobreak(\kongwei)}

\lxopt{A．\(A(x-x_0)+B(y-y_0)=0\)}

\lxopt{B．\(A(x-x_0)-B(y-y_0)=0\)}

\lxopt{C．\(B(x-x_0)+A(y-y_0)=0\)}

\lxopt{D．\(B(x-x_0)-A(y-y_0)=0\)}

\begin{ansblock}[2章-练-课时05-E9]
% ans:2章-练-课时05-E9
\ansitem{9}{A}
\ansnote{详解}{\(M\)在直线上：\(Ax_0+By_0+C=0\)，故\(C=-Ax_0-By_0\)，代入\(Ax+By+C=0\)得\(A(x-x_0)+B(y-y_0)=0\)．选：A．}
\end{ansblock}

\liB{变式\textbf{10}}{简单(知识点三)}{}{已知直线\(l\)：\((a-1)x+(b+2)y+c=0\)，若\(l\parallel y\)轴，但不重合，则下列结论正确的是\nobreak(\kongwei)}

\lxopt{A．\(a\neq1\)，\(b\neq2\)，\(c\neq0\)}

\lxopt{B．\(a\neq1\)，\(b=-2\)，\(c\neq0\)}

\lxopt{C．\(a=1\)，\(b\neq-2\)，\(c\neq0\)}

\lxopt{D．\(a\neq1\)，\(b\neq-2\)，\(c\neq0\)}

\begin{ansblock}[2章-练-课时05-E10]
% ans:2章-练-课时05-E10
\ansitem{10}{B}
\ansnote{详解}{\(l\parallel y\)轴\(\iff\)\(y\)的系数为0且\(x\)的系数不为0：\(b+2=0\)且\(a-1\neq0\)，即\(b=-2\)、\(a\neq1\)；\par\noindent \(l\)与\(y\)轴（方程\(x=0\)）不重合\(\iff c\neq0\)（\(b+2=0\)时\(l\)为\((a-1)x+c=0\)即\(x=-c/(a-1)\)，与\(y\)轴重合\(\iff c=0\)）．\par\noindent 所以\(a\neq1\)，\(b=-2\)，\(c\neq0\)．选：B．}
\end{ansblock}

\liB{变式\textbf{11}}{中档(知识点三)}{}{方程\(y-1=k(x+1)\)在\(k\)取遍所有实数时，可对应无数条直线，这无数条直线的共同点是\nobreak(\kongwei)}

\lxopt{A．都经过定点\((-1,1)\)}

\lxopt{B．都经过定点\((1,-1)\)}

\lxopt{C．斜率都为正}

\lxopt{D．都与\(x\)轴平行}

{\ansblockgrayfalse
\begin{ansblock}[2章-练-课时05-E17]
% ans:2章-练-课时05-E17
\ansitem{11}{A（它们都经过定点(−1,1)；且除直线x＝−1外，过(−1,1)的直线都在其中）}
\ansnote{详解}{对任意实数\(k\)，把\(x=-1\)、\(y=1\)代入方程两边：左\(=1-1=0\)，右\(=k(-1+1)=0\)恒成立，故每条直线都过点\((-1,1)\)；\par\noindent 反之，过\((-1,1)\)且斜率存在（设为\(k\)）的直线方程即\(y-1=k(x+1)\)；斜率不存在的直线\(x=-1\)无法由该方程表示．选：A．（注：本选项组为本卷补写——源题为问答题，为入单选槽按原题数据配四选项，结论为A项不变．）}
\end{ansblock}}

\liB{变式\textbf{12}}{中档(知识点三)}{}{不论\(m\)为何实数，直线\(l\)：\((m-1)x+(2m-1)y=m-5\)恒过一定点，则该定点为\nobreak(\kongwei)}

\bindopt {\fontsize{10.5pt}{\qpTJgLeadLiti}\selectfont \makebox[\dimexpr0.25\linewidth-0.5em\relax][l]{A．\((9,-4)\)}\makebox[\dimexpr0.25\linewidth-0.5em\relax][l]{B．\((-9,4)\)}\makebox[\dimexpr0.25\linewidth-0.5em\relax][l]{C．\((4,-9)\)}\makebox[\dimexpr0.25\linewidth-0.5em\relax][l]{D．\((-4,9)\)}\par}

{\ansblockgrayfalse
\begin{ansblock}[2章-练-课时05-E18]
% ans:2章-练-课时05-E18
\ansitem{12}{A（定点(9,−4)，证明见详解）}
\ansnote{详解}{将方程按\(m\)整理：\(m(x+2y-1)+(-x-y+5)=0\)对一切实数\(m\)恒成立，\par\noindent 须\(x+2y-1=0\)且\(-x-y+5=0\)，联立解得\(x=9\)，\(y=-4\)．\par\noindent 即点\((9,-4)\)使方程对任意\(m\)都成立，故直线\(l\)恒过定点\((9,-4)\)．\par\noindent （取\(m=1\)得\(y=-4\)、取\(m=1/2\)得\(x=9\)，交点\((9,-4)\)代回原方程恒成立，亦可验证．）选：A．（注：本选项组为本卷补写——源题为「求证恒过定点并求坐标」解答形，为入单选槽按原题数据配四选项，定点结论为A项不变，证明详解全保留．）}
\end{ansblock}}

\xiaojie{①识别：以系数条件判定直线位置、由直线上的点反代定式、含参直线系恒过定点为目标的题，判归本题型；一般式与斜截式互化是通法.}
\bindp {\kaishu ②操作：读系数先化斜截式（\(B\neq0\)时\(k=-A/B\)、\(b=-C/B\)），\(B=0\)单列竖直线讨论；含参恒过定点按参数整理，令参数系数与其余因式同时为零解方程组.}
\bindp {\kaishu ③收束：形式互化须防丢解——除以含未知数的式子前要验其是否为零；“平行”与“重合”靠常数项区分，判定后回代检验.}


% ---------- 探究点三 直线方程五种形式的互化与截距式应用 ----------
\tjdnr{三}{直线方程五种形式的互化与截距式应用}{例\textbf{13}}{中档(知识点二)}{已知菱形的两条对角线分别在\(x\)轴和\(y\)轴上，并且它们的长分别等于8和6，求菱形各边所在直线的方程．}
\xiexwei{12mm}

{\ansblockgrayfalse
\begin{ansblock}[2章-练-课时05-E21]
% ans:2章-练-课时05-E21
\ansitem{13}{3x＋4y−12＝0，3x−4y＋12＝0，3x−4y−12＝0，3x＋4y＋12＝0（即3x±4y±12＝0的四种符号组合）}
\ansnote{详解}{对角线在坐标轴上且长分别为8、6，则四个顶点为\((\pm4,0)\)、\((0,\pm3)\)．\par\noindent 四边即依次连接\((4,0)\)、\((0,3)\)、\((-4,0)\)、\((0,-3)\)的回字斜方形，各用截距式：\par\noindent \((4,0)\)-\((0,3)\)：\(x/4+y/3=1\)\(\rightarrow\)\(3x+4y-12=0\)；\((0,3)\)-\((-4,0)\)：\(x/(-4)+y/3=1\)\(\rightarrow\)\(3x-4y+12=0\)；\par\noindent \((-4,0)\)-\((0,-3)\)：\(x/(-4)+y/(-3)=1\)\(\rightarrow\)\(3x+4y+12=0\)；\((0,-3)\)-\((4,0)\)：\(x/4+y/(-3)=1\)\(\rightarrow\)\(3x-4y-12=0\)．}
\end{ansblock}}

\liB{变式\textbf{14}}{中档(知识点二)}{}{已知直线\(l\)过点\(P(2,-1)\)．}

\bindp (1) 若直线\(l\)与直线\(2x+y+3=0\)垂直，求直线\(l\)的方程；
\bindp (2) 若直线\(l\)在两坐标轴的截距互为相反数，求直线\(l\)的方程．
\xiexwei{12mm}

{\ansblockgrayfalse
\begin{ansblock}[2章-练-课时05-E22]
% ans:2章-练-课时05-E22
\ansitem{14}{(1) x−2y−4＝0；(2) x＋2y＝0或x−y−3＝0}
\ansnote{详解}{(1) 与\(2x+y+3=0\)（法向量\((2,1)\)）垂直的直线，其法向量与该直线方向共线，可设\(l\)：\(x-2y+m=0\)（即以\((1,-2)\)为法向量，\((1,-2)\perp(2,1)\)……由课时03点法式知识合法）；\par\noindent 代入\(P(2,-1)\)：\(2+2+m=0\)，\(m=-4\)，故\(l\)：\(x-2y-4=0\)．\par\noindent (2) 当\(l\)过原点时，设\(y=kx\)，代\(P(2,-1)\)得\(k=-1/2\)，即\(x+2y=0\)，此时两截距都是0，互为相反数；\par\noindent 当\(l\)不过原点时，设两截距为\(a\)、\(-a\)（\(a\neq0\)），\(x/a+y/(-a)=1\)，即\(x-y=a\)，代\(P(2,-1)\)：\(a=3\)，\(l\)：\(x-y-3=0\)．\par\noindent 综上，所求直线为\(x+2y=0\)或\(x-y-3=0\)．}
\end{ansblock}}

\liB{变式\textbf{15}}{中档(知识点二)}{}{过点\(P(1,4)\)作直线\(l\)，直线\(l\)与\(x\)、\(y\)轴的正半轴分别交于\(A\)，\(B\)两点，\(O\)为原点．}

\bindp (1) 若\(\triangle ABO\)的面积为9，求直线\(l\)的方程；(2) 若\(\triangle ABO\)的面积为\(S\)，求\(S\)的最小值，并求出此时直线\(l\)的方程．
\xiexwei{14mm}

{\ansblockgrayfalse
\begin{ansblock}[2章-拓-课时05-T6]
% ans:2章-拓-课时05-T6
\ansitem{15}{(1) 2x＋y−6＝0或8x＋y−12＝0；(2) 8，4x＋y−8＝0}
\ansnote{详解}{设\(A(a,0)\)、\(B(0,b)\)，\(a>0\)，\(b>0\)，\(l\)：\(x/a+y/b=1\)，代入\(P(1,4)\)得\(1/a+4/b=1\)（*）．\par\noindent (1) 又\((1/2)ab=9\)即\(ab=18\)；由(*)去分母\(b+4a=ab=18\)，\(b=18-4a\)，\(a(18-4a)=18\)，\(2a^{2}-9a+9=0\)，\(a=3\)或\(3/2\)；\par\noindent 对应\(b=6\)或12，\(l\)为\(x/3+y/6=1\)或\(x/(3/2)+y/12=1\)，即\(2x+y-6=0\)或\(8x+y-12=0\)．\par\noindent (2) \(S=(1/2)ab=(1/2)ab\cdot(1/a+4/b)^{2}=(1/2)(8+b/a+16a/b)\geqslant(1/2)(8+2\sqrt{(b/a)\cdot(16a/b)})=8\)，\par\noindent 当且仅当\(b/a=16a/b\)即\(b=4a\)，代(*)：\(1/a+1/a=1\)，\(a=2\)、\(b=8\)时取等，\(S\)最小值8，此时\(l\)：\(x/2+y/8=1\)即\(4x+y-8=0\)．}
\end{ansblock}}

\liB{变式\textbf{16}}{中档(知识点二)}{}{过点\(P(2,1)\)作直线\(l\)分别交\(x\)轴、\(y\)轴的正半轴于\(A\)，\(B\)两点．}

\bindp (1) 当\(|OA|\cdot|OB|\)取最小值时，求出最小值及直线\(l\)的截距式方程；
\bindp (2) 当\(|PA|\cdot|PB|\)取最小值时，求出最小值及直线\(l\)的截距式方程．
\xiexwei{14mm}

{\ansblockgrayfalse
\begin{ansblock}[2章-拓-课时05-T7]
% ans:2章-拓-课时05-T7
\ansitem{16}{(1) 最小值8，x/4＋y/2＝1；(2) 最小值4，x/3＋y/3＝1}
\ansnote{详解}{(1) 设\(l\)：\(x/a+y/b=1\)（\(a>0\)，\(b>0\)），过\(P\)得\(2/a+1/b=1\)．\par\noindent \(1=2/a+1/b\geqslant2\sqrt{2/(ab)}\)，得\(ab\geqslant8\)，当且仅当\(2/a=1/b\)即\(a=4\)、\(b=2\)时取等；\par\noindent \(|OA|\cdot|OB|=ab\)的最小值为8，此时\(l\)：\(x/4+y/2=1\)．\par\noindent (2) 由(1)得\(b=a/(a-2)\)（\(a>2\)）．\(A(a,0)\)、\(B(0,b)\)，\par\noindent \(|PA|=\sqrt{(a-2)^{2}+1}\)，\(|PB|=\sqrt{4+(b-1)^{2}}\)，而\(b-1=a/(a-2)-1=2/(a-2)\)，\par\noindent \(|PA|\cdot|PB|=\sqrt{(a-2)^{2}+1}\cdot\sqrt{4+4/(a-2)^{2}}\)\par\noindent\(=2\sqrt{(a-2)^{2}+1/(a-2)^{2}+2}\geqslant2\sqrt{2+2}=4\)，\par\noindent 当且仅当\((a-2)^{2}=1/(a-2)^{2}\)即\(a-2=1\)、\(a=3\)时取等，此时\(b=3\)；\par\noindent \(|PA|\cdot|PB|\)的最小值为4，此时\(l\)：\(x/3+y/3=1\)．}
\end{ansblock}}

\liB{变式\textbf{17}}{中档(知识点二)}{}{求分别满足下列条件的直线\(l\)的方程：}

\bindp (1) 斜率是\(3/4\)，且与两坐标轴围成的三角形的面积是6；
\bindp (2) 经过两点\(A(1,0)\)，\(B(m,1)\)；
\bindp (3) 经过点\((4,-3)\)，且在两坐标轴上的截距的绝对值相等．
\xiexwei{14mm}

{\ansblockgrayfalse
\begin{ansblock}[2章-拓-课时05-T8]
% ans:2章-拓-课时05-T8
\ansitem{17}{(1) y＝(3/4)x±3；(2) m≠1时y＝(1/(m−1))(x−1)，m＝1时x＝1；(3) x＋y−1＝0或x−y−7＝0或y＝−(3/4)x}
\ansnote{详解}{(1) 设\(l\)：\(y=(3/4)x+b\)，则\(y\)轴截距\(b\)，\(x\)轴截距\(-4b/3\)，面积\((1/2)|b\cdot(-4b/3)|=2b^{2}/3=6\)，\(b^{2}=9\)，\(b=\pm3\)，\(l\)：\(y=(3/4)x\pm3\)；回代验证：\(b=3\) 代回，面积\((1/2)\cdot|3\cdot(-4)|=6\)；\(b=-3\) 代回，面积\((1/2)\cdot|(-3)\cdot4|=6\)，双向验证成立；\par\noindent (2) \(m\neq1\)时两点横、纵坐标均不等，斜率\(k=(1-0)/(m-1)=1/(m-1)\)，点斜式\(y=(1/(m-1))(x-1)\)；\(m=1\)时\(A\)、\(B\)同在直线\(x=1\)上，\(l\)：\(x=1\)；\par\noindent (3) 设截距\(a\)、\(b\)．①过原点（\(a=b=0\)）：\(l\)过\((0,0)\)、\((4,-3)\)，\(y=-(3/4)x\)；\par\noindent ②不过原点，\(x/a+y/b=1\)，\(|a|=|b|\neq0\)且\(4/a-3/b=1\)：\(a=b\)时\(1/a=1\)、\(a=1\)（\(x+y-1=0\)）；\(a=-b\)时\(4/a+3/a=1\)、\(a=7\)（\(x-y-7=0\)）．\par\noindent 共3条．\par\noindent （按源卷核对注：源答案(1)写作“\(y=(4/3)x\pm3\)”，与其题面“斜率\(3/4\)”不符——本卷按题面取修正值 \(y=(3/4)x\pm3\)；面积回代双向验证已并入(1)详解正文，详见亲算账④“源误更正”．）}
\end{ansblock}}

\xiaojie{①识别：以“求各边方程、按面积与截距条件反求方程、形式互化”为目标的题，判归本题型；两截距齐备先设截距式，含原点疑情单列.}
\bindp {\kaishu ②操作：与两轴围成面积用\(S=|ab|/2\)；最值问题以\(1/a+|x_0|/b=1\)凑基本不等式并验取等；斜率、截距条件先定形式再解参数.}
\bindp {\kaishu ③收束：截距式丢过原点与平行于轴的直线，须补检验；源卷结果与题面互斥时按“源错必改”回代验．}

% ---------- 探究点四 拓展·收难 直线系与包络（需选学§2.2.4） ----------
\tjdnr{四}{拓展·收难 直线系与包络（需选学）}{例\textbf{18}}{难(知识点三)}{（提示：本题详解使用§2.2.4点到直线的距离公式，供已提前选学该条线的学生使用；未学时可只研究题面结构并跳读详解．）}

\bindp 在平面直角坐标系\(xOy\)中，设满足条件\((x-1)\cos\theta+y\sin\theta=1\)（\(0\leqslant\theta<2\pi\)）的所有直线构成集合\(M\)，对于下列四个命题：
\bindp ①\(M\)中所有直线经过一个定点；②存在定点\(P\)不在\(M\)中的任意一条直线上；③\(M\)中的直线所能围成的正三角形的面积都相等；④存在正六边形，使其所在边均在\(M\)中的直线上．其中真命题的序号是\kongbai{}．
\xiexwei{10mm}

{\ansblockgrayfalse
\begin{ansblock}[2章-练-课时05-E19]
% ans:2章-练-课时05-E19
\ansitem{18}{②④}
\ansnote{详解}{根据点到直线距离公式可得点\((1,0)\)到直线\((x-1)\cos\theta+y\sin\theta=1\)的距离\par\noindent \(d=|0\times\cos\theta+0\times\sin\theta-1|/\sqrt{\cos^{2}\theta+\sin^{2}\theta}=1\)，即点\((1,0)\)到直线的距离为定值1，\par\noindent 因此直线\((x-1)\cos\theta+y\sin\theta=1\)是圆\((x-1)^{2}+y^{2}=1\)的切线（\(\theta\)取遍\([0,2\pi)\)恰给出该圆全体切线）．\par\noindent 对于①，\(M\)中所有直线到点\((1,0)\)的距离为定值1，但并不满足经过一个定点，错误；\par\noindent 对于②，当定点\(P\)在该圆内部（不含圆周）时，\(P\)不在\(M\)中的任意一条直线上，正确；\par\noindent 对于③，\(M\)中的直线所能围成的正三角形，既可以是该圆的内接方向的外切正三角形（每边都与圆相切、圆心即内心），也可以让圆恰为三角形的旁切圆型位置（三边仍为切线但圆在三角形外一侧），二者面积不等，错误；\par\noindent 对于④，当正六边形为该圆的外切正六边形时，其六条边所在直线均在\(M\)中，符合要求，正确．\par\noindent 综上所述，真命题的序号是②④．}
\end{ansblock}}

\liB{变式\textbf{19}}{难(知识点三)}{}{（提示：同05-E19，用到点到直线距离公式．）}

\bindp 直线系\(A\)：\((x-3)\cos\alpha+y\sin\alpha=2\)，直线系\(A\)中能组成正三角形的面积等于\kongbai{}．
\xiexwei{10mm}

{\ansblockgrayfalse
\begin{ansblock}[2章-练-课时05-E20]
% ans:2章-练-课时05-E20
\ansitem{19}{4√3/3或12√3}
\ansnote{详解}{直线系\(A\)：\((x-3)\cos\alpha+y\sin\alpha=2\)可变形为\(\sqrt{(x-3)^{2}+y^{2}}\cdot\sin(\alpha+\varphi)=2\)（辅助角公式，\(\tan\varphi=(x-3)/y\)），\par\noindent 故\(\sin(\alpha+\varphi)=2/\sqrt{(x-3)^{2}+y^{2}}\)，而\(|\sin(\alpha+\varphi)|\leqslant1\)，得\((x-3)^{2}+y^{2}\geqslant4\)，\par\noindent 其几何意义为：直线系\(A\)是圆\((x-3)^{2}+y^{2}=4\)的切线的包络（圆上所有点的切线集合）．\par\noindent 把圆心平移至原点，圆\(x^{2}+y^{2}=r^{2}\)上一点\(P(x_0,y_0)\)的切线为\(x_0x+y_0y=r^{2}\)；令\(x_0=r\cos\alpha\)、\(y_0=r\sin\alpha\)，切线即\(x\cos\alpha+y\sin\alpha=r\)，\par\noindent 圆心\((0,0)\)到此直线距离恰为\(r\)，\par\noindent （\(|-r|/\sqrt{\cos^{2}\alpha+\sin^{2}\alpha}=r\)）故\(\alpha\in[0,2\pi)\)时该直线系为圆的全体切线．\par\noindent 直线系中的直线构成正三角形有无数个，但面积的值只有两个．取\(r=2\)：\par\noindent 设一边所在切线为\(y=2\)（上切点水平切线），另两边取斜率\(\pm\sqrt{3}\)的切线\(y=\sqrt{3}x+b\)，圆心到其距离\(|b|/\sqrt{3+1}=2\)，\(b=\pm4\)．\par\noindent ①取\(b\)使三边围成含圆心的正三角形：由\(y=2\)与\(y=\sqrt{3}x-4\)、\(y=-\sqrt{3}x-4\)（即顶点在下方）围成，此时底边\(y=2\)、\(B(-2/\sqrt{3},2)\)，边长\(|BC|=4/\sqrt{3}\)，\(S=(\sqrt{3}/4)\times(4/\sqrt{3})^{2}=4\sqrt{3}/3\)；\par\noindent ②取\(y=-2\)为底、两斜切线为\(y=\sqrt{3}x+4\)、\(y=-\sqrt{3}x+4\)围成另一正三角形：\(B'(-6/\sqrt{3},-2)\)，边长\(|B'C'|=12/\sqrt{3}=4\sqrt{3}\)，\(S=(\sqrt{3}/4)\times(12/\sqrt{3})^{2}=12\sqrt{3}\)．\par\noindent 将圆\(x^{2}+y^{2}=4\)向右平移3个单位即为\((x-3)^{2}+y^{2}=4\)，平移不改变正三角形面积，\par\noindent 故直线系\(A\)中能组成的正三角形面积为\(4\sqrt{3}/3\)或\(12\sqrt{3}\)．}
\end{ansblock}}

\xiaojie{①识别：以直线系命题真假、包络与切线族、面积定值为目标的题，判归本题型；先译出直线系的几何意义（到定点距离为定值＝圆的切线族）.}
\bindp {\kaishu ②操作：含参直线系用“距离定值”或“整理定系数”两法定形；正多边形面积先取标准圆（圆心在原点）算，再平移还原.}
\bindp {\kaishu ③收束：包络型须辨“过定点”与“与定圆相切”两类直线系，结论完全不同；多解（内切、外切）逐一取面积，勿漏第二值.}


% ============================================================
% 课堂评价 G5（恒5＝3单选+2填空＝导槽 G1~G5；ketang 新制顺排可跨栏，禁 \vbox 装栏）
% 印面号 20—24；多选门：本片正文多选 0（题面库多选数0）≤2 合规
% ============================================================
%PROBE ketang块整砍


% ---- 尾块（\tailfill 置 \end{multicols} 前；无参＝内容尾随框，件侧不擅自定高） ----
\tailfill

\end{multicols}
\end{document}
